% TOP_PROBLEM: {"id": 3100029, "title": "What are the (homogeneous adjacency) spectra of the ultracube and the complete hypergraph", "category_id": 2, "category": {"id": 2, "name": "combinatorics", "display_name": "Combinatorics", "description": "Counting problems, graph theory, discrete structures.", "slug": "combinatorics", "order_index": 2, "created_at": "2026-07-31T15:26:25.671Z"}, "status": "open"}
% FIRST_POSTED: null
% ENTRY_KIND: statement_only
% Imported from: batch10.tex; UnsolvedMath ID 3100029
\documentclass[11pt]{article}
\usepackage[margin=25mm]{geometry}
\usepackage{fontspec}
\setmainfont{Latin Modern Roman}
\usepackage{amsmath,amssymb,mathtools}
\usepackage{unicode-math}
\setmathfont{Latin Modern Math}
\usepackage{xurl,hyperref}
\hypersetup{colorlinks=true,urlcolor=blue}
\setcounter{secnumdepth}{0}
\setlength{\parindent}{0pt}
\setlength{\parskip}{5pt}
\setlength{\emergencystretch}{3em}
\newcommand{\sref}[2]{\hyperlink{source:#1:#2}{[S#2]}}
\newcommand{\cataloguescope}[1]{\paragraph{Recorded target.}#1}
\begin{document}
% UnsolvedMath ID 3100029; source catalogue code AMR-030-0029
\setcounter{section}{3100028}
\section{What are the (homogeneous adjacency) spectra of the ultracube and the complete hypergraph}\label{3100029}
{\small\sffamily UnsolvedMath ID 3100029 · Combinatorics}
\subsection{Definitions and mathematical statement}

\paragraph{Shared notation.}$\mathbb N=\{1,2,\ldots\}$, and $\mathbb Z,\mathbb Q,\mathbb R,\mathbb C$ denote the integers, rationals, reals and complex numbers. Any other conventions are specified locally.


For a finite $k$-uniform hypergraph $H=(V,E)$, with $k\geq2$ and $V=[n]$, define its homogeneous adjacency eigenvalues as those $\lambda\in\mathbb C$ for which some $x\in\mathbb C^n\setminus\{0\}$ satisfies
\[\sum_{e\in E:\,i\in e}\prod_{j\in e\setminus\{i\}}x_j=\lambda x_i^{k-1}\quad(i\in[n]).\]
This is the adjacency tensor with entry $1/(k-1)!$ on ordered indices forming an edge, and zero otherwise. The characteristic polynomial is the homogeneous resultant of these $n$ equations after moving the right-hand sides to the left, normalized to be monic in $\lambda$. The spectrum consists of its complex roots, counted with their algebraic multiplicities. The complete $k$-uniform hypergraph $K_n^{(k)}$ has every $k$-subset of $[n]$ as an edge. Let $E_k$ be the hypergraph on $[k]$ consisting of one edge. The $d$-dimensional $k$-uniform ultracube $Q_k^d$ has vertex set $[k]^d$; its edges are the $k$ vertices obtained by varying one coordinate over $[k]$ while fixing all others. Thus $Q_k^d$ is the Cartesian power of $E_k$.
\paragraph{Statement recorded in the supplied source.}
Determine the complete homogeneous adjacency spectra of $K_n^{(k)}$ for all integers $n\geq k\geq2$ and of $Q_k^d$ for all $k\geq2,d\geq1$, including algebraic multiplicities. In particular, describe the ultracube eigenvalues which do not arise merely by summing eigenvalues of its one-edge factors. \sref{3100029}{2}
\subsection{Short English statement}
Determine all homogeneous adjacency eigenvalues, with multiplicity, of complete uniform hypergraphs and Cartesian powers of a single hyperedge.
\subsection{Sources}
\begin{description}
\item[\hypertarget{source:3100029:1}{S1}] ULAM.AI and the UnsolvedMath Contributors, UnsolvedMath: A Curated Collection of Open Mathematics Problems, record 3100029 (AMR-030-0029). CC BY 4.0. \url{https://huggingface.co/datasets/ulamai/UnsolvedMath}.
\item[\hypertarget{source:3100029:2}{S2}] Joshua Cooper and Aaron Dutle, Spectra of Uniform Hypergraphs (2011), Definitions 2.3--2.6 and 3.1; sections 4.3--4.4 and 4.6; section 5, questions (1) and (3). \url{https://arxiv.org/abs/1106.4856}.
\end{description}
\label{end:3100029}
\end{document}
